\subsection{The historical arithmetic development, with its omitted proofs}

\begin{hbdefinition}{def:p12-har-prime-regime}{The exact historical elimination boundary}
The following development reconstructs the chapter's Coq~6.3 example,
not a silently imported modern natural-number library. Its data declaration
is \(\mathsf{nat}:\mathsf{Set}\), with \(O:\mathsf{nat}\) and
\(S:\mathsf{nat}\to\mathsf{nat}\). The licensed operations are
structural recursion and case analysis on this \emph{data} declaration,
and the generated induction principle
\[
 \mathsf{nat\_ind}:\Pi P:\mathsf{nat}\to\mathsf{Prop}.
 P O\to(\Pi n:\mathsf{nat}.P n\to P(S n))\to
 \Pi n:\mathsf{nat}.P n.
\]
The historical logical declarations are likewise retained:
\(\mathsf{and}\ A\ B:\mathsf{Prop}\) has constructor
\(\mathsf{conj}:A\to B\to\mathsf{and}\ A\ B\);
\(\mathsf{TT}:\mathsf{Prop}\) has \(\mathsf{trivial}:\mathsf{TT}\);
\(\mathsf{FF}:\mathsf{Prop}\) has no constructors;
\(\mathsf{or}\ A\ B:\mathsf{Prop}\) has
\(\mathsf{or\_introl}:A\to\mathsf{or}\ A\ B\) and
\(\mathsf{or\_intror}:B\to\mathsf{or}\ A\ B\).
For \(A:\mathsf{Set}\) and \(P:A\to\mathsf{Prop}\),
\(\mathsf{ex}\ A\ P:\mathsf{Prop}\) has
\(\mathsf{ex\_intro}:\Pi x:A.P x\to\mathsf{ex}\ A\ P\).
Here \(\neg a:=a\to\mathsf{FF}\) and
\(a\leftrightarrow b:=(a\to b)\wedge(b\to a)\).
In particular, the displayed source \(\mathsf{Is\_suc}\) case analysis
returns a proposition code, whereas plus, times, predecessor, subtraction
and factorial return data. Induction on \(\mathsf{leseq}\) and elimination
of \(\mathsf{and},\mathsf{or},\mathsf{ex},\mathsf{FF}\) below have
\(\mathsf{Prop}\) targets only. For
\(\mathsf{eq}\ A\ x:\!A\to\mathsf{Prop}\), the constructor is
\(\mathsf{refl\_equal}:\mathsf{eq}\ A\ x\ x\), and its generated
\emph{propositional} eliminator transports \(P x\) along \(x=y\) to
\(P y\), for \(P:A\to\mathsf{Prop}\). No equality reflection,
proof irrelevance, choice, or elimination of an existential proof into
\(\mathsf{Set}\) is used. Finally, retain the explicit source axiom
\[
 \mathsf{DN}:\Pi a:\mathsf{Prop}.((a\to\mathsf{FF})\to\mathsf{FF})\to a.
\]
These are the historical inductive/recursive extensions of the example,
not assertions about the beta-only \(\lambda\mathsf{HOL}\) calculus.
The chapter's historical \((x:A)B\) notation denotes a dependent
product and \([x:A]m\) denotes annotated abstraction; the formulas
below use \(\Pi\) and \(\lambda\) for those same constructions.
Section~3.8's abstract \(\mathsf{nat}:\mathsf{Type}\) and its separately
registered data/proof \(\mathsf{Rec}_{\mathsf{nat}}\) schemas are not
identified with this historical \(\mathsf{Set}\) spelling. The following
proofs are relative to the stated generated principles; a general
consistency theorem for all universe/inductive extensions is a separate
metatheoretic obligation \cite[pp.~1214--1220]{BarendregtGeuvers2001}.
\end{hbdefinition}

\begin{hbdefinition}{def:p12-har-prime-arithmetic}{Source definitions and the needed finite data recursions}
Use the source names \(\mathsf{one}:=S O\), \(\mathsf{two}:=S\mathsf{one}\),
and \(\mathsf{Is\_suc}(O):=\mathsf{FF}\),
\(\mathsf{Is\_suc}(S n):=\mathsf{TT}\). The source equations are
\[
\begin{aligned}
 \mathsf{plus}\ O\ m&=m,& \mathsf{plus}\ (S n)\ m&=S(\mathsf{plus}\ n\ m),\\
 \mathsf{times}\ O\ m&=O,& \mathsf{times}\ (S n)\ m&=
               \mathsf{plus}\ (\mathsf{times}\ n\ m)\ m.
\end{aligned}
\]
Write \(n+m:=\mathsf{plus}\ n\ m\), \(nm:=\mathsf{times}\ n\ m\),
\(n\le m:=\mathsf{leseq}\ n\ m\), and
\(n<m:=\mathsf{leseq}\ (S n)\ m\). The constructors of
\(\mathsf{leseq}\ n\) are \(\mathsf{leseq\_n}:n\le n\) and
\(\mathsf{leseq\_suc}:\Pi m.(n\le m)\to(n\le S m)\).
The instantiated propositional induction schema used for the ordering is
\[
 \Pi n:\mathsf{nat}.\Pi P:\mathsf{nat}\to\mathsf{Prop}.
 P n\to(\Pi m:\mathsf{nat}.(n\le m)\to P m\to P(S m))\to
 \Pi m:\mathsf{nat}.(n\le m)\to P m.
\]
It is generated from these two constructors; it neither changes the
fixed lower bound n nor eliminates into data.
Also \(\mathsf{before}\ n\ P:=\Pi k.(k<n)\to P k\).
The source divisor convention is crucial:
\[
\begin{aligned}
 \mathsf{div}\ d\ n&:=\exists x:\mathsf{nat}.\;xd=n,\\
 \mathsf{propdiv}\ d\ n&:=(\mathsf{one}<d)\wedge(d<n)\wedge\mathsf{div}\ d\ n,\\
 \mathsf{prime}\ n&:=(\mathsf{one}<n)\wedge
                    \neg\exists d.\mathsf{propdiv}\ d\ n,\\
 \mathsf{primediv}\ p\ n&:=\mathsf{prime}\ p\wedge\mathsf{div}\ p\ n,\\
 \mathsf{HPD}\ n&:=\exists p.\mathsf{primediv}\ p\ n.
\end{aligned}
\]
Here \(\exists\) is the source \(\mathsf{ex}:\mathsf{Prop}\), not a
strong sum. Needed data recursions, supplied rather than assumed as
theorems, are
\[
\begin{aligned}
 \mathsf{pred}\ O&:=O,&\mathsf{pred}(S n)&:=n,\\
 O\mathbin{\dot{-}}b&:=O,&(S a)\mathbin{\dot{-}}O&:=S a,&
 (S a)\mathbin{\dot{-}}(S b)&:=a\mathbin{\dot{-}}b,\\
 \mathsf{fac}\ O&:=\mathsf{one},&\mathsf{fac}(S n)&:=(S n)(\mathsf{fac}\ n).
\end{aligned}
\]
Subtraction is structural recursion on the first input, returning a
function of the second; factorial is structural recursion on its input.
All displayed equations are computation equations, not arithmetic
lemmas. More explicitly, the historical \(\mathsf{nat\_rec}\) schema
is instantiated at \(C=\mathsf{nat}\to\mathsf{nat}\) for
\[
\begin{aligned}
 \mathsf{plus}&:=\mathsf{nat\_rec}
  (\lambda m:\mathsf{nat}.m)
  (\lambda n:\mathsf{nat}.\lambda r:\mathsf{nat}\to\mathsf{nat}.
                    \lambda m:\mathsf{nat}.S(r m)),\\
 \mathsf{times}&:=\mathsf{nat\_rec}
  (\lambda m:\mathsf{nat}.O)
  (\lambda n:\mathsf{nat}.\lambda r:\mathsf{nat}\to\mathsf{nat}.
                    \lambda m:\mathsf{nat}.\mathsf{plus}(r m)m),\\
 \mathsf{sub}&:=\mathsf{nat\_rec}
  (\lambda b:\mathsf{nat}.O)
  (\lambda a:\mathsf{nat}.\lambda r:\mathsf{nat}\to\mathsf{nat}.
    \lambda b:\mathsf{nat}.
     \mathsf{nat\_rec}(S a)
       (\lambda c:\mathsf{nat}.\lambda z:\mathsf{nat}.r c)\ b).
\end{aligned}
\]
Here \(a\mathbin{\dot{-}}b:=\mathsf{sub}\ a\ b\), and the inner
sub recursor has \(C=\mathsf{nat}\) and ignores its recursive z.
At that same data target,
\(\mathsf{pred}:=\mathsf{nat\_rec}\ O\
(\lambda n:\mathsf{nat}.\lambda r:\mathsf{nat}.n)\)
and \(\mathsf{fac}:=\mathsf{nat\_rec}\ \mathsf{one}\
(\lambda n:\mathsf{nat}.\lambda r:\mathsf{nat}.\mathsf{times}(S n)r)\).
The base and step of each instance have respectively types C and
\(\mathsf{nat}\to C\to C\); their constructor computation gives
exactly the equations above. This uses several instantiated recursion
schemas, not a new base-calculus constant polymorphic over a top sort.
The source's repeated \((n,m,k:\mathsf{nat})\) binder groups in
\(\mathsf{plus\_ass}\) and \(\mathsf{distr}\) are reconstructed below
with one binder group: no hidden extra parameters are required
\cite[pp.~1219--1220 and p.~1215, note~4]{BarendregtGeuvers2001}.
If the repeated names are instead read as literal binder shadowing,
alpha-rename the outer triple to \(n_0,m_0,k_0\). The resulting
six-argument type has those three arguments unused; its witness is
\(\lambda n_0:\mathsf{nat}.\lambda m_0:\mathsf{nat}.
\lambda k_0:\mathsf{nat}.L\), where L is the fully quantified
single-triple proof of plus_ass or distr constructed below. Thus neither
reading requires an additional arithmetic lemma; no claim about the
historical parser accepting that typography is needed.
\end{hbdefinition}

\begin{hblemma}{thm:p12-har-prime-logical}{proved=annex:p12.har.prime-logical}{The source logical library, with classical dependence separated}
Let \(a,b,c:\mathsf{Prop}\), \(A:\mathsf{Set}\), and
\(P,Q:A\to\mathsf{Prop}\); all these parameters are universally
quantified in the following signatures.
The source constructor/eliminator aliases
\(\mathsf{and\_in},\mathsf{and\_ell},\mathsf{and\_elr},
\mathsf{false\_el},\mathsf{or\_inl},\mathsf{or\_inr},\mathsf{or\_el}\)
have the generated propositional introduction/projection/elimination
types. Also \(\mathsf{all\_el}:\Pi x:A.(\Pi y:A.P y)\to P x\)
and \(\mathsf{ex\_in}:\Pi x:A.P x\to\mathsf{ex}\ A\ P\)
are constructed directly, as is
\(\mathsf{non\_or}:\neg(a\vee b)\to\neg a\wedge\neg b\).
The source \(\mathsf{non\_ex},\mathsf{all\_not},\mathsf{all\_and},
\mathsf{ex\_or}\) have the following types:
\[
\begin{aligned}
 \mathsf{non\_ex}&:\neg(\exists x.P x)\to\Pi x.\neg P x,\\
 \mathsf{all\_not}&:(\Pi x.\neg P x)\to\neg\exists x.P x,\\
 \mathsf{all\_and}&:(\Pi x.P x\wedge Q x)\to
                         (\Pi x.P x)\wedge(\Pi x.Q x),\\
 \mathsf{ex\_or}&:(\exists x.P x\vee Q x)\to
                         (\exists x.P x)\vee(\exists x.Q x).
\end{aligned}
\]
Each bound x has domain A. The source
\(\mathsf{dn\_c}:\neg\neg a\leftrightarrow a\),
\(\mathsf{tnd}:a\vee\neg a\),
\(\mathsf{non\_all}:\neg(\forall x.P x)\to\exists x.\neg P x\),
and \(\mathsf{ex\_c}:(\exists x.P x)\leftrightarrow\neg\forall x.\neg P x\)
are constructed using the displayed \(\mathsf{DN}\). The first four
displayed predicate-logic lemmas and the introduction/elimination aliases
need no classical axiom and require no inhabitant of \(A\).
\end{hblemma}

\begin{hblemma}{thm:p12-har-prime-equality}{proved=annex:p12.har.prime-equality}{Equality transport and natural-number separation}
The source \(\mathsf{sym\_eq},\mathsf{leib},\mathsf{eq\_ind\_r},
\mathsf{f\_equal},\mathsf{no\_conf}\) are derivable from the stated
propositional equality eliminator, at these exact signatures:
\[
\begin{aligned}
 \mathsf{sym\_eq}&:\Pi A:\mathsf{Set}.\Pi x,y:A.(x=y)\to(y=x),\\
 \mathsf{leib}&:\Pi A:\mathsf{Set}.\Pi P:A\to\mathsf{Prop}.
                   \Pi x,y:A.(x=y)\to P x\to P y,\\
 \mathsf{eq\_ind\_r}&:\Pi A:\mathsf{Set}.\Pi x:A.
    \Pi P:A\to\mathsf{Prop}.P x\to\Pi y:A.(y=x)\to P y,\\
 \mathsf{f\_equal}&:\Pi A,B:\mathsf{Set}.\Pi f:A\to B.
                        \Pi x,y:A.(x=y)\to(f x=f y),\\
 \mathsf{no\_conf}&:\Pi n:\mathsf{nat}.\neg(O=S n).
\end{aligned}
\]
Moreover,
\(S n=S m\to n=m\). None uses \(\mathsf{DN}\).
\end{hblemma}

\begin{hblemma}{thm:p12-har-prime-order}{proved=annex:p12.har.prime-order}{The inductive ordering and course-of-values induction}
The source \(\mathsf{leseq\_trans},\mathsf{lthan\_leseq},
\mathsf{non\_lt0},\mathsf{suc\_leseq},\mathsf{lt01},
\mathsf{n0n1lt},\mathsf{cv\_ind}\) are derivable:
\[
\begin{aligned}
 \mathsf{leseq\_trans}&:\Pi x,y,z.(x\le y)\to(y\le z)\to(x\le z),\\
 \mathsf{lthan\_leseq}&:\Pi n,m.(n<m)\to(n\le m),\\
 \mathsf{non\_lt0}&:\Pi n.\neg(n<O),\\
 \mathsf{suc\_leseq}&:\Pi n,m.(S n\le S m)\to(n\le m),\\
 \mathsf{lt01}&:\Pi x.(x=O)\vee(x=\mathsf{one})\vee(\mathsf{one}<x),\\
 \mathsf{n0n1lt}&:\Pi n.\neg(n=O)\to\neg(n=\mathsf{one})\to(\mathsf{one}<n).
\end{aligned}
\]
The omitted natural-number domains are all \(\mathsf{nat}\). Useful derived
forms are zero minimality, trichotomy, antisymmetry, strict
irreflexivity, \(k<S n\leftrightarrow k\le n\), and
\(k\le S n\to(k\le n\vee k=S n)\). In particular,
\[
 (\Pi n.\mathsf{before}\ n\ P\to P n)\to\Pi n.P n
\]
is proved, not supplied as a new induction axiom.
\end{hblemma}

\begin{hblemma}{thm:p12-har-prime-semiring}{proved=annex:p12.har.prime-semiring}{The actual addition and multiplication library}
The source \(\mathsf{plus\_altsuc},\mathsf{plus\_altzero},
\mathsf{plus\_ass},\mathsf{plus\_com},\mathsf{distr},
\mathsf{timesaltzero},\mathsf{timesaltsuc},\mathsf{times\_ass}\)
are proved at the following single-binder types (all variables in nat):
\[
\begin{aligned}
 \mathsf{plus\_altsuc}&:\Pi n,m.n+S m=S(n+m),\\
 \mathsf{plus\_altzero}&:\Pi n.n+O=n,\\
 \mathsf{plus\_ass}&:\Pi n,m,k.n+(m+k)=(n+m)+k,\\
 \mathsf{plus\_com}&:\Pi n,m.n+m=m+n,\\
 \mathsf{distr}&:\Pi n,m,k.(n+m)k=nk+mk,\\
 \mathsf{timesaltzero}&:\Pi n.nO=O,\\
 \mathsf{timesaltsuc}&:\Pi n,m.n(S m)=nm+n,\\
 \mathsf{times\_ass}&:\Pi n,m,k.n(mk)=(nm)k.
\end{aligned}
\]
In addition,
\(nm=mn\), \(\mathsf{one}\,n=n\,\mathsf{one}=n\), both addition
cancellation laws, and
\[
 n\le m\quad\leftrightarrow\quad\exists k.\;n+k=m
\]
are derived. For \(O<d\), multiplication by \(d\) preserves strict
order and reflects weak order. These facts require no \(\mathsf{DN}\).
\end{hblemma}

\begin{hblemma}{thm:p12-har-prime-divisibility}{proved=annex:p12.har.prime-divisibility}{The subtraction dependency actually discharged}
Divisibility is reflexive and transitive, and \(\mathsf{one}\) divides
every natural number. A number \(d>\mathsf{one}\) cannot divide
\(\mathsf{one}\). For every \(d,a,b\), including \(d=O\),
\[
 \mathsf{div}\ d\ a\to\mathsf{div}\ d\ (a+b)\to\mathsf{div}\ d\ b.
\]
For local coefficients \(rd=a\), \(sd=a+b\) and \(O<d\), the
constructed coefficient of \(b\) is \(s\mathbin{\dot{-}}r\).
This is a proof of a \(\mathsf{Prop}\) existential, not a function
extracting a witness from an arbitrary existential proof into
\(\mathsf{Set}\).
\end{hblemma}

\begin{hbtheorem}{thm:p12-har-prime-divisor}{proved=annex:p12.har.prime-divisor}{The source prime-divisor theorem, now with its proof}
The historical declaration has an actual derivation:
\[
 \mathsf{numbers\_gt1\_have\_primediv}:\Pi n:\mathsf{nat}.
             (\mathsf{one}<n)\to\mathsf{HPD}\ n.
\]
Its only classical step is \(\mathsf{tnd}\) on
\(\exists d.\mathsf{propdiv}\ d\ n\), derived from
\(\mathsf{DN}\); the remaining steps use course-of-values induction
and the proved transitivity of divisibility.
\end{hbtheorem}

\begin{hbtheorem}{thm:p12-har-prime-euclid}{proved=annex:p12.har.prime-euclid}{The factorial continuation and Euclid's unboundedness statement}
For every \(N:\mathsf{nat}\),
\[
 \exists p:\mathsf{nat}.\;\mathsf{prime}\ p\wedge N<p.
\]
The proof supplies the factorial bound and the theorem that every
\(\mathsf{one}\le m\le N\) divides \(\mathsf{fac}\ N\), then
applies the source prime-divisor theorem to \(\mathsf{fac}\ N+\mathsf{one}\).
It inherits that theorem's explicit \(\mathsf{DN}\) dependence.
No constructive prime-producing program follows by unlicensed
\(\mathsf{Prop}\)-to-\(\mathsf{Set}\) elimination. This is the first
formula of the source footnote. The footnote's informal variable-length
list display is a different representation; identifying it with a single
coded first-order arithmetic formula requires an explicit coding bridge,
not an assertion that first-order arithmetic cannot express finite lists
\cite[p.~1215, note~4; p.~1220]{BarendregtGeuvers2001}.
There is also a fully typed finite-family formulation within the same
regime. For any \(k:\mathsf{nat}\) and \(a:\mathsf{nat}\to\mathsf{nat}\),
\[
 \exists p.\;\mathsf{prime}\ p\wedge
           \Pi i:\mathsf{nat}.(i<k)\to\neg(p=a i).
\]
The proof bounds the first k values by the actually constructed data
fold \(B O\ a:=O\), \(B(S k)\ a:=B k\ a+a k\), then uses the
unboundedness result. It does not require the listed values to be prime,
so it entails the source's finite-prime-family avoidance assertion.
This dependent-function representation supplies no unproved
identification with an arithmetic sequence-coding library.
\end{hbtheorem}
